\chapter{Universal class and finite-stage representability}
\label{route-ch-universal}

Use the carrier construction and comparison of
Chapter~\ref{route-ch-finite-stage}.  All chart, class, and cocycle data will
be defined over one finite Galois subextension \(L/K\) of \(\Omega/K\).
The resulting package \(P\) has \(L=P.N.1\) and carrier
\(J_L=P.\mathrm{gluedOver}\).  This same carrier is used throughout.

For a field extension \(E'/E\) and an \(E'\)-scheme \(T\), we use the
canonical identification between \(\operatorname{Pic}^0_{C_E}(T)\) and
\(\operatorname{Pic}^0_{C_{E'}}(T)\).  Associativity of fibre products gives
the identification on invertible sheaves and on classes modulo pullbacks
from \(T\); it respects etale sheafification and field-valued degree.
These identifications commute with pullbacks and with further field
extension.

\section{What is pinned over the separably closed field}

\begin{theorem}[The universal class and all atlas restrictions over \(\Omega\)]
  \label{route-thm-sepclosed-universal-class}
  \lean{AlgebraicGeometry.pic0SepClosedUniversalClass,
    AlgebraicGeometry.pic0FiniteStageUniversalChartClass,
    AlgebraicGeometry.pic0FiniteStageUniversalOverlapClass,
    AlgebraicGeometry.pic0FiniteStageUniversalChartClass_restrict_left,
    AlgebraicGeometry.pic0FiniteStageUniversalChartClass_restrict_right,
    AlgebraicGeometry.Pic0FiniteStageUniversalAtlasClass,
    AlgebraicGeometry.pic0FiniteStageUniversalAtlasClass}
  \leanok
  \dcref{custom: Yoneda universal element and its restrictions}
  \uses{route-thm-sepclosed-representability, route-thm-finite-atlas}

  Let
  \[
    \xi_\Omega:=r_\Omega(\mathrm{id}_{J_\Omega})
       \in\operatorname{Pic}^0_{C_\Omega}(J_\Omega).
  \]
  For every selected affine chart \(U_i\subset J_\Omega\) and every exact
  ordered pair-overlap \(U_{ij}=U_i\cap U_j\), the checked classes are exactly
  the pullbacks of this one \(\xi_\Omega\).  If
  \(q^i_{ij}:U_{ij}\to U_i\) and \(q^j_{ij}:U_{ij}\to U_j\), then
  \[
    (q^i_{ij})^*(\xi_\Omega|_{U_i})
      =\xi_\Omega|_{U_{ij}}
      =(q^j_{ij})^*(\xi_\Omega|_{U_j}).
  \]
\end{theorem}

\begin{proof}
  Apply the representing equivalence to the identity of \(J_\Omega\).  Its
  naturality says that precomposition by a chart or overlap morphism is
  pullback of the resulting Picard class, proving both restriction equations.
  Every class in this theorem still lives over \(\Omega\).
\end{proof}

\section{Reflection of equality}

\begin{lemma}[Faithfully flat reflection for relative Picard classes]
  \label{route-lem-relative-picard-reflection}
  \lean{AlgebraicGeometry.relPicAlgMap_faithfullyFlat_injective}
  \dcref{custom: relative Picard faithful-flat kernel lemma}
  \uses{route-conv-objects}

  Let \(E\) be a field extension of \(K\) and let \(A\to B\) be a
  faithfully flat map of \(E\)-algebras.  Pullback is injective on
  \[
    \operatorname{Pic}(C_A)/\operatorname{im}\operatorname{Pic}(\Spec A)
      \longrightarrow
    \operatorname{Pic}(C_B)/\operatorname{im}\operatorname{Pic}(\Spec B).
  \]
\end{lemma}

\begin{proof}
  Write \(p_A:C_A\to\Spec A\).  Properness and geometric integrality of
  \(C_E\) give \(H^0(C_E,\mathcal O)=E\), and hence
  \((p_A)_*\mathcal O_{C_A}=\mathcal O_{\Spec A}\) for every \(A\).
  Indeed a finite affine cover of the separated curve computes global
  sections by an equalizer, and tensoring over the field \(E\) preserves
  that equalizer.  Thus this assertion also holds over \(B\).

  Suppose an invertible sheaf \(\mathcal M\) on \(C_A\) becomes
  \(p_B^*\mathcal N\) for an invertible sheaf \(\mathcal N\) on
  \(\Spec B\).  Flat base change for global sections and the projection
  formula identify the pullback of \((p_A)_*\mathcal M\) with
  \((p_B)_*p_B^*\mathcal N=\mathcal N\).  Faithfully flat descent of
  finite locally free modules implies that \((p_A)_*\mathcal M\) is
  invertible.  The evaluation map
  \(p_A^*(p_A)_*\mathcal M\to\mathcal M\) becomes an isomorphism over
  \(B\), so it is an isomorphism over \(A\).  The original relative
  class is therefore zero.  Apply this kernel statement to the difference
  of two classes.
\end{proof}

\begin{lemma}[Faithfully flat reflection on affine degree-zero classes]
  \label{route-lem-pic0-affine-reflection}
  \lean{AlgebraicGeometry.pic0Map_faithfullyFlat_injective}
  \dcref{custom: affine etale plus construction and relative Picard reflection}
  \uses{route-lem-relative-picard-reflection}

  For the same \(E\) and \(A\to B\), the map
  \(\operatorname{Pic}^0_{C_E}(\Spec A)\to
    \operatorname{Pic}^0_{C_E}(\Spec B)\) is injective.
\end{lemma}

\begin{proof}
  The preceding lemma applied to faithfully flat etale covers shows that
  the relative Picard presheaf is separated on affine tests.  Its map to
  its etale sheafification is consequently injective.  Represent a
  sheafified class \(z\) over \(A\) by a relative Picard class \(x\)
  on an affine etale cover \(\Spec R\to\Spec A\).
  If \(z\) becomes zero over \(B\), then over \(R\otimes_A B\) its
  restriction is the image of the restriction of \(x\).
  Injectivity of the relative-to-sheafified map makes that relative
  restriction zero.  Since \(R\to R\otimes_A B\) is faithfully flat,
  Lemma~\ref{route-lem-relative-picard-reflection} gives \(x=0\).
  Etale separation now gives \(z=0\).  This proves injectivity on the
  sheafified Picard group, and hence on its degree-zero subgroup.
\end{proof}

\begin{lemma}[Affine faithfully flat reflection on arbitrary tests]
  \label{route-lem-pic0-field-reflection}
  \lean{AlgebraicGeometry.pic0Map_affine_faithfullyFlat_injective,
    AlgebraicGeometry.pic0CrossBaseMap_injective}
  \leanok
  \dcref{custom: affine faithful-flat reflection and Zariski separation}
  \uses{route-lem-pic0-affine-reflection}

  Let \(E\) be a field, let \(D/E\) be a smooth proper geometrically
  irreducible curve, and let \(f:T'\to T\) be an affine faithfully
  flat morphism of \(E\)-schemes. Then pullback
  \[
    f^*:\operatorname{Pic}^0_D(T)\longrightarrow
      \operatorname{Pic}^0_D(T')
  \]
  is injective. In particular, for every field extension \(E'/E\),
  simultaneous scalar extension of the curve and test gives an injection
  \[
    \operatorname{Pic}^0_D(T)\longrightarrow
       \operatorname{Pic}^0_{D_{E'}}(T\times_E E').
  \]
  No quasi-compactness assumption on \(T\) is needed.
\end{lemma}

\begin{proof}
  Cover \(T\) by affine opens \(U_i=\Spec A_i\). Since \(f\) is
  affine, each inverse image is affine, say
  \(V_i=f^{-1}(U_i)=\Spec B_i\). The restriction \(V_i\to U_i\) is
  flat and surjective; therefore its map on section rings
  \(A_i\to B_i\) is faithfully flat. If \(f^*x=f^*y\), the
  restrictions of \(x\) and \(y\) to \(U_i\) agree after pullback
  to \(V_i\). Lemma~\ref{route-lem-pic0-affine-reflection} gives
  equality on every \(U_i\), and Zariski separation of the etale
  Picard sheaf gives \(x=y\) on \(T\).

  For the field-extension assertion, the projection
  \(p:T\times_E E'\to T\) is affine faithfully flat, as the base
  change of \(\Spec E'\to\Spec E\). Apply the first assertion to
  \(p\), then compose with the canonical scalar-extension equivalence
  \(\operatorname{Pic}^0_D(T\times_E E')\cong
    \operatorname{Pic}^0_{D_{E'}}(T\times_E E')\).
\end{proof}

\section{The joint finite Galois stage}

\begin{theorem}[Joint Galois-stage package and descended universal class]
  \label{route-thm-descended-universal-class}
  \notready
  \source{stacks-project}
  \dcref{Tags 01ZM, 0BMG, and 0EXM; Tag 0B9K as a source boundary}
  \uses{route-thm-sepclosed-universal-class, route-thm-global-comparison,
    route-lem-pic0-affine-reflection}

  There exist, simultaneously:
  \begin{enumerate}
  \item a legacy package
    \(P:\operatorname{Pic0FiniteStageGluePackage}(C_\Omega,K)\);
  \item for \(L:=P.N.1\subset\Omega\), certificates that \(L/K\) is finite
    Galois, with the displayed inclusion and scalar-tower maps;
  \item on the literal carrier \(J_L:=P.\mathrm{gluedOver}\), a class
    \[
      \xi_L\in\operatorname{Pic}^0_{C_L}(J_L)
        =\operatorname{pic0Subgroup}
          ((\operatorname{baseChange}_{K,L})C,P.\mathrm{gluedOver}),
      \qquad C_L=C\times_KL.
    \]
  \end{enumerate}
  The symbol \(P\) in this theorem is a new, jointly selected package; it is
  not the arbitrary comparison witness \(P_0\) from
  Theorem~\ref{route-thm-global-comparison}.  If the construction starts from
  \(P_0\), it must additionally provide an explicit refinement map from
  \(P_0\) to \(P\) together with the induced \(\operatorname{Over}\)-map and
  its compatibility with every chart and overlap comparison.  Without those
  fields, \(P_0\) is discarded after the comparison theorem and no field
  enlargement is allowed to silently replace its carrier.  All subsequent
  occurrences of \(J_L\), \(\xi_L\), (FS-Rep), (QP), and (OA) quantify over
  this single final \(P\).
  For the comparison
  \(\beta_P:J_L\times_L\Omega\xrightarrow{\sim}J_\Omega\), this class must
  satisfy
  \[
    \operatorname{bc}_{L,\Omega}(\xi_L)=\beta_P^*(\xi_\Omega)
      \quad\text{in}\quad
    \operatorname{Pic}^0_{C_\Omega}(J_L\times_L\Omega).
    \tag{UC}
  \]
  Equivalently, its restriction to every descended chart pulls back to
  \(\xi_\Omega|_{U_i}\), and its two restrictions to every descended overlap
  agree before gluing.
\end{theorem}

\begin{proof}
  Start with finite-presentation models of the chart and overlap diagram.
  On each chart the restriction of \(\xi_\Omega\) is etale-locally a
  relative invertible-sheaf class.  Quasi-compactness permits finitely many
  affine etale charts of finite presentation.  Invertible sheaves, their
  gluing isomorphisms, and witnesses of equality modulo classes from the
  base have finite presentations on these charts.  The etale descent
  equations are witnessed on finitely many affine refinements of their
  overlaps.  Applying finite-presentation descent to these objects,
  morphisms, and equations, together with the carrier diagram, defines
  the finitely many sheafified chart classes over a common finite stage.
  This argument uses the actual etale equivalence relation, rather than
  an identification with a different Picard sheaf.

  Enlarge that stage to its finite normal closure \(L\) inside \(\Omega\).
  Since \(\Omega/K\) is separable, \(L/K\) is finite Galois.  Base change
  all the descended data to \(L\), and assemble the package \(P\) there.
  In particular the chart classes and their restriction maps are on
  the charts of this \(P\).  Their two restrictions to each affine
  overlap agree after extension to \(\Omega\), because both come from
  \(\xi_\Omega\).  Lemma~\ref{route-lem-pic0-affine-reflection}, applied
  first to the ambient sheafified Picard groups as in its proof, reflects
  those equalities to \(L\).

  Each descended class has degree zero.  Indeed, given any field-valued
  point with field \(F\) on a descended chart, the nonzero ring
  \(F\otimes_L\Omega\) has a residue field \(F'\).  The induced point
  over \(\Omega\) has degree zero, and invariance of degree under
  \(F\to F'\) gives degree zero at the original point.
  The compatible chart classes thus glue uniquely in the degree-zero
  etale sheaf to \(\xi_L\) on \(J_L\).  After base change to
  \(\Omega\), its restrictions on all the charts agree with those of
  \(\beta_P^*\xi_\Omega\), by construction of the chart comparisons.
  Zariski separation proves (UC).  No global representative by an
  invertible sheaf on \(C_L\times_L J_L\) is asserted.
\end{proof}

\section{Representability from a compatible universal class}

\begin{theorem}[Descent of representability along a field extension]
  \label{route-thm-field-descent-representability}
  \lean{AlgebraicGeometry.pic0RepresentableBy_of_fieldExtension,
    AlgebraicGeometry.pic0RepresentableBy_of_fieldExtension_homEquiv}
  \leanok
  \source{stacks-project}
  \dcref{Tag 040L, descent of morphisms}
  \uses{route-thm-sepclosed-representability,
    route-lem-pic0-field-reflection}

  Let \(E\) be a field, let \(\Omega/E\) be a field extension with
  \(\Omega\) separably closed, and let \(D/E\) be a smooth proper
  geometrically irreducible curve. Let \((R_\Omega,r_\Omega)\) be the
  representing pair for \(\operatorname{Pic}^0_{D_\Omega}\), and put
  \(\eta_\Omega=r_\Omega(\mathrm{id}_{R_\Omega})\). Suppose that an
  \(E\)-scheme \(J\), a class \(\xi\in\operatorname{Pic}^0_D(J)\),
  and an \(\Omega\)-isomorphism \(\beta:J_\Omega\to R_\Omega\)
  satisfy
  \[
    \operatorname{bc}_{E,\Omega}(\xi)=\beta^*\eta_\Omega
      \quad\text{in}\quad\operatorname{Pic}^0_{D_\Omega}(J_\Omega).
  \]
  Then, for every \(E\)-scheme \(T\), pullback gives a bijection
  \[
    \Phi_T:\Hom_E(T,J)\xrightarrow{\sim}\operatorname{Pic}^0_D(T),
    \qquad f\longmapsto f^*\xi.
  \]
  These bijections are natural in \(T\), so \(J\) represents
  \(\operatorname{Pic}^0_D\) with universal class \(\xi\).
\end{theorem}

\begin{proof}
  Pullback of \(\xi\) defines \(\Phi_T\); functoriality of pullback
  proves naturality. For every \(\Omega\)-scheme \(S\), the
  base-change adjunction, \(\beta\), and \(r_\Omega\) give a
  bijection
  \[
    \Hom_E(S,J)\simeq\Hom_\Omega(S,J_\Omega)
      \simeq\operatorname{Pic}^0_{D_\Omega}(S)
      \simeq\operatorname{Pic}^0_D(S).
  \]
  Naturality and the assumed equality of universal classes identify
  its value on each morphism with pullback of \(\xi\).

  Let \(p:T_\Omega\to T\), and let
  \(a\in\operatorname{Pic}^0_D(T)\). The preceding bijection supplies
  a unique \(E\)-morphism \(f_0:T_\Omega\to J\) with
  \(f_0^*\xi=p^*a\). Put \(W=T_\Omega\times_TT_\Omega\), with
  projections \(p_1,p_2\), and give \(W\) its first
  \(\Omega\)-structure. Both \(f_0p_1\) and \(f_0p_2\) are
  \(E\)-morphisms. Their pullbacks of \(\xi\) agree because
  \(pp_1=pp_2\). The displayed bijection for the
  \(\Omega\)-scheme \(W\) gives \(f_0p_1=f_0p_2\). The
  base-change adjunction applies to both \(E\)-morphisms, regardless
  of whether \(p_2\) preserves the chosen \(\Omega\)-structure.

  The morphism \(p\) is affine faithfully flat. Fpqc descent of
  scheme morphisms, \cite[Tag 040L]{stacks-project}, gives a unique
  \(f:T\to J\) with \(fp=f_0\). The equality
  \(p^*(f^*\xi)=p^*a\) is reflected by
  Lemma~\ref{route-lem-pic0-field-reflection}, so \(\Phi_T(f)=a\).
  If \(\Phi_T(f)=\Phi_T(g)\), the bijection on \(T_\Omega\)
  implies \(fp=gp\), and uniqueness in morphism descent gives
  \(f=g\). Thus \(\Phi_T\) is bijective for every \(T\).
  The identity of \(J\) pulls \(\xi\) back to itself, proving the
  assertion about the universal class. Neither \(T\) nor the
  extension \(\Omega/E\) is required to satisfy a finiteness
  assumption.
\end{proof}

\begin{theorem}[Yoneda equivalence on the same finite Galois carrier]
  \label{route-thm-finite-stage-yoneda}
  \notready
  \uses{route-thm-descended-universal-class,
    route-thm-field-descent-representability}

  For the same \((P,L,\xi_L)\), every \(L\)-scheme \(T\) admits the
  bijection
  \[
    \Phi_T:\Hom_L(T,J_L)\xrightarrow{\sim}
       \operatorname{Pic}^0_{C_L}(T),
    \qquad f\longmapsto f^*\xi_L.
    \tag{Y}
  \]
  For every \(q:T'\to T\) and \(f:T\to J_L\), the naturality equation is
  \[
    \Phi_{T'}(f\circ q)=q^*\Phi_T(f).
  \]
  Consequently \(J_L=P.\mathrm{gluedOver}\) represents
  \(\operatorname{Pic}^0_{C_L}\):
  \[
    h_{J_L}\xrightarrow{\sim}\operatorname{Pic}^0_{C_L}.
    \tag{FS-Rep}
  \]
\end{theorem}

\begin{proof}
  Use the canonical associativity isomorphism
  \((C\times_K L)\times_L\Omega\cong C\times_K\Omega\) to
  identify their degree-zero Picard functors. Compose \(\beta_P\)
  with the unique isomorphism between the corresponding representing
  schemes. Naturality of that isomorphism identifies their universal
  classes, so (UC) is the universal-class compatibility required by
  Theorem~\ref{route-thm-field-descent-representability}. Apply that
  theorem with \(E=L\), \(D=C_L\), \(J=P.\mathrm{gluedOver}\),
  and \(\xi=\xi_L\). Its evaluation formula and naturality give
  (Y) and (FS-Rep).
\end{proof}

\begin{remark}[The universal class determines the transformation]
  \label{route-rem-first-frontier}
  \dcref{custom}

  For a fixed \(J_L\), Yoneda identifies natural transformations
  \(h_{J_L}\to\operatorname{Pic}^0_{C_L}\) with classes on \(J_L\).
  It does not assert that such a transformation is invertible.  Equation
  (UC) and the reflection and descent arguments above prove invertibility
  for the class constructed here.
\end{remark}
